\section{Esports and Shoutcasting}
\label{sec:introduction:esports_and_shoutcasting}

Electronic sports, often shortened to esports, have experienced growth as part of the entertainment industry of videogames, with significant media coverage and economic impact \cite{block2021esports}. As a result of this growth, there has been a professionalization of the broadcasting of esports events, with higher production levels, enhancing the viewer experience for the audiences. In contrast to traditional sports, which heavily rely on physical presence, the need for audience engagement in esports is necessary for industry growth through the formation of online communities and fanbases \cite{koronios2020100373}.

\medskip

To facilitate the audience engagement, esports take a similar approach to traditional sports. In both professional and amateur esports events, ``shoutcasters'' are often featured, and are responsible for providing live commentary and analysis of the videogame gameplay. Depending on their role, shoutcasters are usually categorized into play-by-play casters, who describe the ongoing action, and color casters, who provide expert analysis and insights \cite{sell2015sports}. Although this two-caster format is common and these two types of discourses are the most frequent in esports events, games with more complex or slower-paced gameplay often give rise to additional discursive functions beyond real-time narration and analysis.

\medskip

For this research, our domain focuses on the shoutcasting of the videogame \textit{League of Legends} (LoL), developed by Riot Games, and released in 2009. LoL is a Multiplayer Online Battle Arena (MOBA) game, where two teams of five players compete against each other to destroy the opposing team's base. To this day, it presents itself as one of the most popular esports games, with a peak of 117 million active players across worldwide servers.

\medskip

The game's nature involves heavy micro-management from each individual player, which also relies on a heavy understanding of the different playable characters, alongside strategic decision-making as a collective effort from the teams \cite{10.1145/3769534.3769592}. Additionally, matches are characterized by alternating phases of high-intensity action and slower strategic gameplay. This makes LoL an ideal case for the study of the discursive functions of shoutcasting. In addition, match durations, which range anywhere from 20 to 60 minutes, allow for a variety of discursive functions to be observed across different extended periods.

\medskip

Despite the heavy requirements of communication strategies, the existing literature regarding the discursive functions and language structure of shoutcasting in esports is still limited. In the following section, we will review some of these roles in greater depth and their relevance to audience engagement in esports events. This will provide a foundation for the modeling of the discursive functions of shoutcasting in LoL.